\subsection{\texorpdfstring{紧集的 \(\mu\) -稠密族}{紧集的 mu -稠密族}}

命题 12. --- 设 \(\mu\) 是局部紧空间 \(X\) 上的一个测度，A 是一个 \(\mu\)-可测的、含于 \(X\) 的子集，\(\mathfrak{K}\) 是 A 的紧子集所成的一个集族，满足下列条件：

\((\mathrm{PL}_{\mathrm{I}})\) 属于 \(\mathfrak{K}\) 的每个集合的闭（从而为紧）子集都属于 \(\mathfrak{K}\) 。\((\mathrm{PL}_{\mathrm{II}})\) \(\mathfrak{K}\) 中集合的每个有限并都属于 \(\mathfrak{K}\) 。

下列四个性质彼此等价：

\begin{enumerate}
\def\labelenumi{\alph{enumi})}
\item
  为使 A 的子集 B 是局部 \(\mu\)-可忽略的，其充要条件是：\(|\mu|^*(\mathrm{B} \cap \mathrm{K}) = 0\) 对所有 \(\mathrm{K} \in \mathfrak{K}\) 成立。
\item
  对 A 的每个紧子集 \(K_{0}\) 与每个 \(\varepsilon > 0\)，存在一个集 \(K \in \mathfrak{K}\)，它含于 \(K_{0}\) 且使 \(|\mu|(\mathrm{K}_{0} - \mathrm{K}) \leqslant \varepsilon\) 。
\item
  对 A 的每个紧子集 B，存在 B 的一个分割，它由一个 \(\mu\)-可忽略集 N 与一个紧集序列 \((\mathrm{H}_{n})\)（属于 \(\mathfrak{K}\)）组成。
\item
  对 A 的每个紧子集 B，存在一个递增的紧集序列 \((\mathrm{K}_{n})\)（各项均属于 \(\mathfrak{K}\)），含于 B，且使集 \(N = B - \bigcup K_{n}\) 为 \(\mu\)-可忽略。
\end{enumerate}

直接可得（No.~2, 命题 5）：\(d\) 蕴涵 \(a\)；即知 \(c\) 蕴涵 \(d\)：取 \(\mathrm{K}_n\) 为诸 \(\mathrm{H}_p\)（\(p \leqslant n\)）之并并引用 \((\mathrm{PL}_{\mathrm{II}})\)。为证 \(b\) 蕴涵 \(c\)，我们递归地定义一个集合序列 \((\mathrm{H}_p)\)（由 \(\mathfrak{K}\) 中的集合组成），使得 \(\mathrm{H}_{n+1} \subset \mathrm{B} - \bigcup_{p \leqslant n} \mathrm{H}_p\) 且 \(|\mu| (\mathrm{B} - \bigcup_{p \leqslant n} \mathrm{H}_p) \leqslant 1/n\)（ \(\S 4\) , No.~6, 定理 4）。

剩下要证 \(a)\) 蕴涵 \(b)\) 。用反证法，设 \(\alpha\)（诸数 \(|\mu|(\mathrm{K})\) 的上确界，其中 K 取遍 \(\mathrm{K}_0\) 的属于 \(\mathfrak{K}\) 的子集）是 \(< |\mu|(\mathrm{K}_0)\) 的。由 \((\mathrm{PL}_{\mathrm{II}})\)，存在一个递增序列 \((\mathrm{L}_n)\)（由 \(\mathrm{K}_0\) 的紧子集组成，属于 \(\mathfrak{K}\)），且 \(\sup_n |\mu|(\mathrm{L}_n) = \alpha\) 。令 \(\mathrm{B} = \bigcup_n \mathrm{L}_n\)；B 可积且 \(|\mu|(\mathrm{B}) = \alpha\)，因此 \(|\mu|(\mathrm{K}_0 - \mathrm{B}) = |\mu|(\mathrm{K}_0) - \alpha > 0\) 。另一方面，我们将看到，对每个集 \(\mathrm{K} \in \mathfrak{K}\) 都有 \(|\mu|(\mathrm{K} \cap (\mathrm{K}_0 - \mathrm{B})) = 0\)，而由 \(a)\)，这将导出矛盾。事实上，若存在集 \(\mathrm{K} \in \mathfrak{K}\) 使 \(|\mu|(\mathrm{K} \cap (\mathrm{K}_0 - \mathrm{B})) > 0\)，则存在 \(\mathrm{K} \cap (\mathrm{K}_0 - \mathrm{B})\) 的一个紧子集 H，使 \(|\mu|(\mathrm{H}) > 0\) 。由 \((\mathrm{PL}_{\mathrm{I}})\)，将有 \(\mathrm{H} \in \mathfrak{K}\)，且当 \(n\) 充分大时，

\[
| \mu | (\mathrm{L} _ {n} \cup \mathrm{H}) = | \mu | (\mathrm{L} _ {n}) + | \mu | (\mathrm{H}) > \alpha .
\]

但 \(L_{n} \cup H\) 属于 \(\mathfrak{K}\)（由 \((\mathrm{PL}_{\mathrm{II}})\)），这与 \(\alpha\) 的定义矛盾。

定义 6. --- 设 A 是 X 的一个 \(\mu\)-可测子集。称 A 的紧子集所成的一个集族 \(\mathfrak{K}\) 在 A 中 \(\mu\)-稠密，如果它满足命题 12 的条件 \((\mathrm{PL}_{\mathrm{I}})\) 、 \((\mathrm{PL}_{\mathrm{II}})\) 、 \(a)\) 、 \(b)\) 、 \(c)\) 、 \(d)\) 。

A 的全体紧子集所成之集族在 A 中 \(\mu\)-稠密。

当 A = X 时，我们简称“ \(\mu\)-稠密集族”以代替“在 X 中 \(\mu\)-稠密的集族”。若 X - A 局部 \(\mu\)-可忽略，则在 A 中 \(\mu\)-稠密的每个紧子集族也在 X 中 \(\mu\)-稠密。

注. --- 设 A 是一个紧集序列 \((\mathrm{L}_{n})\) 与一个 \(\mu\)-可忽略（相应地，局部 \(\mu\)-可忽略）集之并，\(\mathfrak{K}\) 是在 A 中 \(\mu\)-稠密的一个紧子集族。对每个 \(L_{n}\) 应用命题 12 陈述中的性质 c)，我们看出，A 是一个由属于 \(\mathfrak{K}\) 的紧集组成的序列与一个 \(\mu\)-可忽略（相应地，局部 \(\mu\)-可忽略）集之并。

若 K 是 X 的一个紧子集，则“K 的一个紧子集族在 K 中 \(\mu\)-稠密”与“它在 K 中 \(\mu_{K}\)-稠密”是一回事；这由 No.~7 的引理 2 与引理 3 以及命题 12 的条件 b) 得出。

命题 13. --- 设 A 是 X 的一个 \(\mu\)-可测子集，\(\mathfrak{K}\) 是在 A 中 \(\mu\)-稠密的一个紧子集族。设 \(\mathfrak{H}\) 是 A 的紧子集所成的一个集族，它满足 (PL \(_{\text{I}}\) ) 与 (PL \(_{\text{II}}\) )，且对每个 K \(\in\) \(\mathfrak{K}\)，满足 H \(\in\) \(\mathfrak{H}\) 与 H \(\subset\) K 的那些 H 所成之集族在 K 中 \(\mu_{\text{K}}\)-稠密（或等价地说，\(\mu\)-稠密）。则 \(\mathfrak{H}\) 在 A 中 \(\mu\)-稠密。

因为，设 L 是 A 的一个紧子集。对每个 \(\varepsilon > 0\)，存在 \(K \in \mathfrak{K}\)，使得 \(K \subset L\) 且 \(|\mu|(\mathrm{L} - \mathrm{K}) \leqslant \varepsilon/2\)；然后存在 \(H \in \mathfrak{H}\)，使得 \(H \subset K\) 且 \(|\mu|(\mathrm{K} - \mathrm{H}) \leqslant \varepsilon/2\)；由此得 \(|\mu|(\mathrm{L} - \mathrm{H}) \leqslant \varepsilon\)，命题得证。

